\section{非极化准夸克分布的单圈结果}

本节考虑非极化准夸克分布 $\tilde q(x, \mu^2, P^z)$ 的单圈修正。由于非阿贝尔性质不进入非单态情形，非单态夸克分布的单圈计算与 QED 类似。

由于单圈结果是规范不变的，我们可以任选规范进行计算。最简单的选择是轴向规范 $A^z=0$~\cite{axial}，因为在此规范中式~(\ref{qUnpolDef})中的规范链接（gauge link）化为单位元。在轴向规范中，相关的费曼图如图 1 所示，其中非局域算符用虚线表示。这些图中包含紫外（UV）发散、软发散和共线发散。我们用夸克质量 $m$ 正规化共线发散。软发散预计在各图之间抵消。紫外发散则用横向动量截断正规化。当然，这种截断破坏转动对称性；然而预计它能给出正确的领头对数行为。

\begin{figure}[htbp]
\includegraphics[width=0.2\textwidth]{images/vertex.pdf}
\hspace{10em}
\includegraphics[width=0.2\textwidth]{images/self.pdf}
\caption{准夸克分布的单圈修正图。}
\label{onelooppdf}
\end{figure}

单圈图给出的结果为
\begin{equation}
                    \tilde q(x, \mu^z, P^z) = (1+\tilde Z^{(1)}_F+\dots) \delta(x-1) + \tilde q^{(1)}(x) + \dots
\end{equation}
其中
\begin{align}\label{unpolvertexnoexp}
\tilde q^{(1)}(x)&=\frac{\alpha_S C_F}{2\pi} \left\{ \begin{array} {ll} \frac{1+x^2}{1-x}\ln \frac{x(\Lambda(x)-x P^z)}{(x-1)(\Lambda(1-x)+P^z(1-x))}+1-\frac{x P^z}{\Lambda(x)}+\frac{x\Lambda(1-x)+(1-x)\Lambda(x)}{(1-x)^2 P^z}\ , & x>1\ , \\
\ \\
\frac{1+x^2}{1-x}\ln\frac{(P^z)^2}{m^2}+\frac{1+x^2}{1-x}\ln \frac{4x(\Lambda(x)-x P^z)}{(1-x)(\Lambda(1-x)+(1-x)P^z)}-\frac{4x}{1-x}+1-\frac{x P^z}{\Lambda(x)}\ & \\
+\frac{x\Lambda(1-x)+(1-x)\Lambda(x)}{(1-x)^2 P^z}\ , & 0<x<1\ , \\
\ \\
\frac{1+x^2}{1-x}\ln \frac{(x-1)(\Lambda(x)-x P^z)}{x(\Lambda(1-x)+(1-x)P^z)}-1-\frac{x P^z}{\Lambda(x)}+\frac{x\Lambda(1-x)+(1-x)\Lambda(x)}{(1-x)^2P^z}\ , & x<0 \end{array} \right.
\end{align}
对有限的 $P^z$ 成立，其中 $\Lambda(x)=\sqrt{\mu^2+x^2 (P^z)^2}$，而含共线发散的对数与标准的 Altarelli-Parisi 劈裂核~\cite{AP}相联系。波函数重正化修正同样依赖于动量 $P^z$：
\begin{align}\label{selfnoexp}
\tilde Z^{(1)}_F&=\frac{\alpha_S C_F}{2\pi} \int dy \left\{ \begin{array} {ll} -\frac{1+y^2}{1-y}\ln \frac{y(\Lambda(y)-y P^z)}{(y-1)(\Lambda(1-y)+P^z(1-y))}-1-\frac{y\Lambda(1-y)+(1-y)\Lambda(y)}{(1-y)^2 P^z}\ & \\
+\frac{y^2 P^z}{\Lambda(y)}+\frac{y(1-y)P^z}{\Lambda(1-y)}+\frac{\Lambda(y)-\Lambda(1-y)}{P^z}\ , & y>1\ , \\
\ \\
-\frac{1+y^2}{1-y}\ln\frac{(P^z)^2}{m^2}-\frac{1+y^2}{1-y}\ln \frac{4y(\Lambda(y)-y P^z)}{(1-y)(\Lambda(1-y)+(1-y)P^z)}+\frac{4y^2}{1-y}+1\ & \\
-\frac{y\Lambda(1-y)+(1-y)\Lambda(y)}{(1-y)^2P^z}+\frac{y^2 P^z}{\Lambda(y)}+\frac{y(1-y)P^z}{\Lambda(1-y)}+\frac{\Lambda(y)-\Lambda(1-y)}{P^z}\ , & 0<y<1\ , \\
\ \\
-\frac{1+y^2}{1-y}\ln \frac{(y-1)(\Lambda(y)-y P^z)}{y(\Lambda(1-y)+(1-y)P^z)}+1-\frac{y\Lambda(1-y)+(1-y)\Lambda(y)}{(1-y)^2P^z}\ & \\
+\frac{y(1-y)P^z}{\Lambda(1-y)}+\frac{y^2 P^z}{\Lambda(y)}+\frac{\Lambda(y)-\Lambda(1-y)}{P^z}\ , & y<0\ . \end{array} \right.
\end{align}
容易验证，该结果在单圈精度上满足矢量流守恒条件
\begin{equation}
  \int^{+\infty}_{-\infty} dx\, \tilde q(x, \mu^2, P^z) =1
\end{equation}
由于准分布中夸克的组分并不具有部分子诠释，部分子动量分数的取值范围从 $-\infty$ 一直延伸到 $+\infty$。有趣的是，共线发散只存在于 $0<x<1$ 区域。


在场论计算中，人们通常把紫外截断取为大于问题中的任何其他标度。换言之，应取 $\mu\rightarrow \infty$ 极限，仅保留领头贡献而忽略幂次压低项。原则上，格点 QCD 计算也应当如此。于是，准分布的实际场论结果应为
\begin{align}\label{unpolvertex}
\tilde q^{(1)}(x,\mu,P^z)&=\frac{\alpha_S C_F}{2\pi} \left\{ \begin{array} {ll} \frac{1+x^2}{1-x}\ln \frac{x}{x-1}+1+\frac{\mu}{(1-x)^2P^z}\ , & x>1\ , \\ \frac{1+x^2}{1-x}\ln \frac{(P^z)^2}{m^2}+\frac{1+x^2}{1-x}\ln\frac{4x}{1-x}-\frac{4x}{1-x}+1+\frac{\mu}{(1-x)^2P^z}\ , & 0<x<1\ , \\ \frac{1+x^2}{1-x}\ln \frac{x-1}{x}-1+\frac{\mu}{(1-x)^2P^z}\ , & x<0\ , \end{array} \right.
\end{align}
以及
\begin{align}\label{unpolself}
\tilde Z^{(1)}_F
=\frac{\alpha_S C_F}{2\pi} \int dy \left\{ \begin{array} {ll} -\frac{1+y^2}{1-y}\ln \frac{y}{y-1}-1-\frac{\mu}{(1-y)^2P^z}\ , & y>1\ , \\ -\frac{1+y^2}{1-y}\ln \frac{(P^z)^2}{m^2}-\frac{1+y^2}{1-y}\ln\frac{4y}{1-y}+\frac{4y^2}{1-y}+1-\frac{\mu}{(1-y)^2P^z}\ , & 0<y<1\ , \\ -\frac{1+y^2}{1-y}\ln \frac{y-1}{y}+1-\frac{\mu}{(1-y)^2P^z}\ , & y<0\ . \end{array} \right.
\end{align}
上述结果有几个有趣的特征值得注意：1）在 $x>1$ 与 $x<0$ 区域也存在贡献。其背后的物理是透明的：当母粒子具有有限动量 $P^z$ 时，组分部分子的动量可以大于 $P^z$，甚至可以为负。这与无限动量系的结果截然不同——后者中动量分数被限制在 $0<x<1$。2）存在一个来自规范链接自能的线性发散；若换到非轴向规范（例如 Feynman 规范），这一点很容易看出。这种线性发散在维数正规化中通常被忽略，但它确实出现在格点正规化中。3）不存在对数型紫外发散。相反，在 $0<x<1$ 区域存在对 $P^z$ 的对数依赖。后面我们将看到，通过一个匹配条件，这种对数依赖可以被转化为光锥部分子分布的重正化标度依赖。4）所有软发散均已抵消；但仍有以夸克质量依赖形式体现的共线发散残留。

另一方面，在同一正规化方案下，可以通过取 $P^z\rightarrow \infty$ 极限来计算光锥部分子分布。这按照文献~\cite{Weinberg:1966jm}的精神完成，结果为
\begin{equation}
                    q(x, \mu^2) = (1+ Z^{(1)}_F+ \dots) \delta(x-1) +  q^{(1)}(x) + \dots
\end{equation}
其中
\begin{align}\label{unpolvertexIMF}
q^{(1)}(x)&=\frac{\alpha_S C_F}{2\pi} \left\{ \begin{array} {ll} 0\ , & x>1\  \mbox{或}\  x<0\ , \\ \frac{1+x^2}{1-x}\ln \frac{\mu^2}{m^2}-\frac{1+x^2}{1-x}\ln{(1-x)^2}-\frac{2x}{1-x}\ , & 0<x<1\ ,\end{array} \right.
\end{align}
以及
\begin{align}\label{unpolselfIMF}
Z^{(1)}_F &=\frac{\alpha_S C_F}{2\pi} \int dy \left\{ \begin{array} {ll} 0 \ , & y>1 \ \mbox{或} \ y<0\ , \\ -\frac{1+y^2}{1-y}\ln \frac{\mu^2}{m^2}+\frac{1+y^2}{1-y}\ln{(1-y)^2}+\frac{2y}{1-y}\ , & 0<y<1\ , \end{array} \right.
\end{align}
其中 $\delta Z^{(1)}$ 的被积函数恰与 $q^{(1)}(x)$ 的相反。该结果与在横向动量截断方案中由部分子分布的光锥定义所导出的结果完全一致。而且，共线（质量）奇异性显然与准部分子分布中的相同。这表明在单圈水平上，准部分子分布捕捉了无限动量系中的全部共线物理。此外，贡献仅来自中间胶子可被切割的那幅图，而该图具有部分子诠释。
